\documentclass[10pt,a4paper]{amsart}
\usepackage{amsmath,amssymb,amsthm,graphicx}
% Retained preamble for the Stacks Project chapter snapshot in this directory.
% The chapter's own preamble.tex is a copy of the shared Stacks Project preamble; the block
% below is sliced MECHANICALLY, line for line, from it. Nothing was added, removed or
% redefined; the wrapper class, the xr-hyper/hyperref wiring and the externaldocument list
% are omitted because the assembler supplies the document class.
%
% lines taken from preamble.tex: 3, 5, 9, 11, 12, 156, 157, 158, ... (51 lines total)
%
% package coverage restored from the chapter's own preamble (2026-10-13): xy, xy-2cell

\usepackage[all]{xy}
\xyoption{2cell}
\UseAllTwocells

\newenvironment{reference}{\comment}{\endcomment}
\newenvironment{slogan}{\comment}{\endcomment}
\newenvironment{history}{\comment}{\endcomment}
\theoremstyle{plain}
\newtheorem{theorem}[subsection]{Theorem}
\newtheorem{proposition}[subsection]{Proposition}
\newtheorem{lemma}[subsection]{Lemma}
\theoremstyle{definition}
\newtheorem{definition}[subsection]{Definition}
\newtheorem{example}[subsection]{Example}
\newtheorem{exercise}[subsection]{Exercise}
\newtheorem{situation}[subsection]{Situation}
\theoremstyle{remark}
\newtheorem{remark}[subsection]{Remark}
\newtheorem{remarks}[subsection]{Remarks}
\numberwithin{equation}{subsection}
\def\lim{\mathop{\mathrm{lim}}\nolimits}
\def\colim{\mathop{\mathrm{colim}}\nolimits}
\def\Spec{\mathop{\mathrm{Spec}}}
\def\Hom{\mathop{\mathrm{Hom}}\nolimits}
\def\Ext{\mathop{\mathrm{Ext}}\nolimits}
\def\SheafHom{\mathop{\mathcal{H}\!\mathit{om}}\nolimits}
\def\SheafExt{\mathop{\mathcal{E}\!\mathit{xt}}\nolimits}
\def\Sch{\mathit{Sch}}
\def\Mor{\mathop{\mathrm{Mor}}\nolimits}
\def\Ob{\mathop{\mathrm{Ob}}\nolimits}
\def\Sh{\mathop{\mathit{Sh}}\nolimits}
\def\NL{\mathop{N\!L}\nolimits}
\def\CH{\mathop{\mathrm{CH}}\nolimits}
\def\proetale{{pro\text{-}\acute{e}tale}}
\def\etale{{\acute{e}tale}}
\def\QCoh{\mathit{QCoh}}
\def\Ker{\mathop{\mathrm{Ker}}}
\def\Im{\mathop{\mathrm{Im}}}
\def\Coker{\mathop{\mathrm{Coker}}}
\def\Coim{\mathop{\mathrm{Coim}}}
\DeclareMathSymbol{\boxtimes}{\mathbin}{AMSa}{"02}
\def\QCohstack{\mathcal{QC}\!\mathit{oh}}
\def\Cohstack{\mathcal{C}\!\mathit{oh}}
\def\Spacesstack{\mathcal{S}\!\mathit{paces}}
\def\Quotfunctor{\mathrm{Quot}}
\def\Hilbfunctor{\mathrm{Hilb}}
\def\Curvesstack{\mathcal{C}\!\mathit{urves}}
\def\Polarizedstack{\mathcal{P}\!\mathit{olarized}}
\def\Complexesstack{\mathcal{C}\!\mathit{omplexes}}
\def\Pic{\mathop{\mathrm{Pic}}\nolimits}
\def\Picardstack{\mathcal{P}\!\mathit{ic}}
\def\Picardfunctor{\mathrm{Pic}}
\def\Deformationcategory{\mathcal{D}\!\mathit{ef}}
% lane-s11: discard bibliographic \begin{reference} blocks entirely (incl. any \cite inside).
\newsavebox{\limrefbox}
\ifcsname reference\endcsname\expandafter\let\csname reference\endcsname\undefined\fi\ifcsname endreference\endcsname\expandafter\let\csname endreference\endcsname\undefined\fi
\def\citation#1{\relax}
\newenvironment{reference}{\setbox\limrefbox=\vbox\bgroup\let\cite\relax}{\egroup\ignorespaces}
\usepackage{cleveref}
\crefname{originalreference}{original reference}{original references}
\Crefname{originalreference}{Original reference}{Original references}
\makeatletter
\newlabel{morphisms-lemma-quasi-affine-finite-type-over-S}{{04II}{source}{}{sourcepointer-1}{}}
\newlabel{morphisms-lemma-quasi-affine-finite-type-over-S@cref}{{[originalreference][0][]04II}{[1][0][]source}{}{}{}}
\makeatother
\crefname{originalreference}{original reference}{original references}
\Crefname{originalreference}{Original reference}{Original references}
\makeatletter
\newlabel{morphisms-lemma-quasi-projective-open-projective}{{07RM}{source}{}{sourcepointer-2}{}}
\newlabel{morphisms-lemma-quasi-projective-open-projective@cref}{{[originalreference][0][]07RM}{[1][0][]source}{}{}{}}
\makeatother
\crefname{originalreference}{original reference}{original references}
\Crefname{originalreference}{Original reference}{Original references}
\makeatletter
\newlabel{morphisms-section-scheme-theoretic-closure}{{01RA}{source}{}{sourcepointer-3}{}}
\newlabel{morphisms-section-scheme-theoretic-closure@cref}{{[originalreference][0][]01RA}{[1][0][]source}{}{}{}}
\makeatother
\crefname{originalreference}{original reference}{original references}
\Crefname{originalreference}{Original reference}{Original references}
\makeatletter
\newlabel{morphisms-lemma-image-proper-is-proper}{{03GN}{source}{}{sourcepointer-4}{}}
\newlabel{morphisms-lemma-image-proper-is-proper@cref}{{[originalreference][0][]03GN}{[1][0][]source}{}{}{}}
\makeatother
\crefname{originallemma}{original lemma}{original lemmas}
\Crefname{originallemma}{Original lemma}{Original lemmas}
\makeatletter
\newlabel{lemma-Noetherian-dvr-valuative-proper}{{0208}{source}{}{sourcepointer-5}{}}
\newlabel{lemma-Noetherian-dvr-valuative-proper@cref}{{[originallemma][0][]0208}{[1][0][]source}{}{}{}}
\makeatother
\crefname{originalreference}{original reference}{original references}
\Crefname{originalreference}{Original reference}{Original references}
\makeatletter
\newlabel{schemes-lemma-separated-implies-valuative}{{01KZ}{source}{}{sourcepointer-6}{}}
\newlabel{schemes-lemma-separated-implies-valuative@cref}{{[originalreference][0][]01KZ}{[1][0][]source}{}{}{}}
\makeatother
\crefname{originalreference}{original reference}{original references}
\Crefname{originalreference}{Original reference}{Original references}
\makeatletter
\newlabel{schemes-lemma-separated-permanence}{{01KU}{source}{}{sourcepointer-7}{}}
\newlabel{schemes-lemma-separated-permanence@cref}{{[originalreference][0][]01KU}{[1][0][]source}{}{}{}}
\makeatother
\crefname{originalreference}{original reference}{original references}
\Crefname{originalreference}{Original reference}{Original references}
\makeatletter
\newlabel{morphisms-lemma-characterize-proper}{{0BX5}{source}{}{sourcepointer-8}{}}
\newlabel{morphisms-lemma-characterize-proper@cref}{{[originalreference][0][]0BX5}{[1][0][]source}{}{}{}}
\makeatother
\title{Refined valuative criterion for properness along limits}
\author{The Stacks Project Authors}
\date{stacks-project:limits.tex:local-snapshot-20261012; GNU FDL 1.2 or later}
\begin{document}
\maketitle
\noindent\textit{Editorial note.} The reduction to the affine valuation ring case is the standard locality ('immediate'); the constructive content is fully worked in the counted span.\par\medskip
\noindent\textit{Allowed standard prerequisites (editorial).}\begin{itemize}
\item Properness of limits. Limits, properness lemmas of the chapter
\item Valuative criterion. Morphisms, valuative criterion lemmas
\end{itemize}\medskip
\ifdefined\hypertarget\hypertarget{sourcepointer-1}{}\fi
\noindent\textit{Source reference.} Original reference 04II: sources/stacks-chow/tags.txt, source lines 5904--5904. Cross-chapter standard prerequisite cited by the author inside the counted span; it is used as a pointer and its own proof is not part of this item.\par\medskip
\ifdefined\hypertarget\hypertarget{sourcepointer-2}{}\fi
\noindent\textit{Source reference.} Original reference 07RM: sources/stacks-chow/tags.txt, source lines 9872--9872. Cross-chapter standard prerequisite cited by the author inside the counted span; it is used as a pointer and its own proof is not part of this item.\par\medskip
\ifdefined\hypertarget\hypertarget{sourcepointer-3}{}\fi
\noindent\textit{Source reference.} Original reference 01RA: sources/stacks-chow/tags.txt, source lines 2295--2295. Cross-chapter standard prerequisite cited by the author inside the counted span; it is used as a pointer and its own proof is not part of this item.\par\medskip
\ifdefined\hypertarget\hypertarget{sourcepointer-4}{}\fi
\noindent\textit{Source reference.} Original reference 03GN: sources/stacks-chow/tags.txt, source lines 4534--4534. Cross-chapter standard prerequisite cited by the author inside the counted span; it is used as a pointer and its own proof is not part of this item.\par\medskip
\ifdefined\hypertarget\hypertarget{sourcepointer-5}{}\fi
\noindent\textit{Source reference.} Original lemma 0208: sources/stacks-chow/tags.txt, source lines 2625--2625. Cross-chapter standard prerequisite cited by the author inside the counted span; it is used as a pointer and its own proof is not part of this item.\par\medskip
\ifdefined\hypertarget\hypertarget{sourcepointer-6}{}\fi
\noindent\textit{Source reference.} Original reference 01KZ: sources/stacks-chow/tags.txt, source lines 2068--2068. Cross-chapter standard prerequisite cited by the author inside the counted span; it is used as a pointer and its own proof is not part of this item.\par\medskip
\ifdefined\hypertarget\hypertarget{sourcepointer-7}{}\fi
\noindent\textit{Source reference.} Original reference 01KU: sources/stacks-chow/tags.txt, source lines 2063--2063. Cross-chapter standard prerequisite cited by the author inside the counted span; it is used as a pointer and its own proof is not part of this item.\par\medskip
\ifdefined\hypertarget\hypertarget{sourcepointer-8}{}\fi
\noindent\textit{Source reference.} Original reference 0BX5: sources/stacks-chow/tags.txt, source lines 14968--14968. Cross-chapter standard prerequisite cited by the author inside the counted span; it is used as a pointer and its own proof is not part of this item.\par\medskip
\begin{lemma}
\label{lemma-refined-valuative-criterion-proper}
Let $f : X \to S$ and $h : U \to X$ be morphisms of schemes.
Assume that $S$ is locally Noetherian, that $f$ and $h$ are of finite type,
that $f$ is separated, and that $h(U)$ is dense in $X$.
If given any commutative solid diagram
$$
\xymatrix{
\Spec(K) \ar[r] \ar[d] & U \ar[r]^h & X \ar[d]^f \\
\Spec(A) \ar[rr] \ar@{-->}[rru] & & S
}
$$
where $A$ is a discrete valuation ring with field of fractions $K$, there
exists a dotted arrow making the diagram commute, then $f$ is proper.
\end{lemma}

\begin{proof}
There is an immediate reduction to the case where $S$ is affine.
Then $U$ is quasi-compact.
Let $U = U_1 \cup \ldots \cup U_n$ be an affine open covering.
We may replace $U$ by $U_1 \amalg \ldots \amalg U_n$ without
changing the assumptions, hence we may assume $U$ is affine.
Thus we can find an open immersion $U \to Y$ over $X$
with $Y$ proper over $X$. (First put $U$ inside $\mathbf{A}^n_X$
using Morphisms, Lemma \ref{morphisms-lemma-quasi-affine-finite-type-over-S}
and then take the closure inside $\mathbf{P}^n_X$, or you can directly
use Morphisms, Lemma \ref{morphisms-lemma-quasi-projective-open-projective}.)
We can assume $U$ is dense in $Y$ (replace $Y$ by the scheme theoretic
closure of $U$ if necessary, see Morphisms, Section
\ref{morphisms-section-scheme-theoretic-closure}).
Note that $g : Y \to X$ is surjective as the image is closed
and contains the dense subset $h(U)$.
We will show that $Y \to S$ is proper. This will imply that
$X \to S$ is proper by
Morphisms, Lemma \ref{morphisms-lemma-image-proper-is-proper}
thereby finishing the proof.
To show that $Y \to S$ is proper we will use
part (4) of Lemma \ref{lemma-Noetherian-dvr-valuative-proper}.
To do this consider a diagram
$$
\xymatrix{
\Spec(K) \ar[r]_y \ar[d] & Y \ar[d]^{f \circ g} \\
\Spec(A) \ar[r] \ar@{..>}[ru] & S
}
$$
where $A$ is a discrete valuation ring with fraction field $K$
and where $y : \Spec(K) \to Y$ is the inclusion of a generic point.
We have to show there exists a unique dotted arrow.
Uniqueness holds by the converse to the valuative criterion
for separatedness
(Schemes, Lemma \ref{schemes-lemma-separated-implies-valuative})
since $Y \to S$ is separated as the
composition of the separated morphisms $Y \to X$ and $X \to S$
(Schemes, Lemma \ref{schemes-lemma-separated-permanence}).
Existence can be seen as follows.
As $y$ is a generic point of $Y$, it is contained in $U$.
By assumption of the lemma
there exists a morphism $a : \Spec(A) \to X$ such that
$$
\xymatrix{
\Spec(K) \ar[r]_y \ar[d] & U \ar[r] & X \ar[d]^f \\
\Spec(A) \ar[rr] \ar[rru]^a & & S
}
$$
is commutative. Then since $Y \to X$ is proper, we can
apply the valuative criterion for properness
(Morphisms, Lemma \ref{morphisms-lemma-characterize-proper})
to find a morphism $b : \Spec(A) \to Y$ such that
$$
\xymatrix{
\Spec(K) \ar[r]_y \ar[d] & Y \ar[d]^g \\
\Spec(A) \ar[r]^a \ar[ru]^b & X
}
$$
is commutative. This finishes the proof since
$b$ can serve as the dotted arrow above.
\end{proof}

\end{document}
